% !TeX root = ../../Book1.tex
\section{弱-*\WeakStarJoin{}拓扑}
\label{b1:sec:0902}

% 来源：Fremlin2016Measure，第2卷附录2A5I(g)，附录第27页，实际阅读作者公开PDF。
% 弱-*与弱拓扑的有限测试解释据此改写；l1对偶、坐标判别与尾部示性序列在此完整推导。
当向量本身就是泛函时，有一种更直接的测试办法：给定一个输入，观察泛函在该输入上的值。这样的测试只使用原空间的向量，而不使用对偶空间上的全部连续线性泛函，因而产生了另一种拓扑。弱拓扑与弱-*\WeakStarJoin{}拓扑的区别，首先是允许使用的测试族不同。

\subsection{对偶空间上的弱-*\WeakStarJoin{}拓扑}
\label{b1:subsec:090201}

\begin{definition}\label{b1:def:weak-star-topology}
设 \(X\) 为赋范空间。对 \(f\in X^*\)、有限集 \(E\subset X\) 与 \(\varepsilon>0\)，令
\[
 V(f;E,\varepsilon)
 \coloneqq\{\,g\in X^*\mid |g(x)-f(x)|<\varepsilon\text{ 对每个 }x\in E\,\}.
\]
以这些集合为邻域基的拓扑称为 \(X^*\) 上的\term[ruo-xing-tuopu]{弱-*\WeakStarJoin{}拓扑}[weak-* topology]，记为 \(\sigma(X^*,X)\)。
\end{definition}
\symidx{\(\sigma(X^*,X)\)：对偶空间相对于原空间的弱-*\WeakStarJoin{}拓扑}

与\cref{b1:prop:weak-topology-properties}中的有限测试论证完全相同，这些集合组成开邻域基，而 \(\sigma(X^*,X)\) 是使所有求值映射
\[
 X^*\longrightarrow\mathbb K,\quad f\longmapsto f(x)\quad(x\in X)
\]
连续的最弱拓扑。若 \(f\ne g\)，总有一个 \(x\in X\) 使 \(f(x)\ne g(x)\)，故两个小求值邻域可将它们分开，弱-*\WeakStarJoin{}拓扑也是 Hausdorff 的。估计 \(|(g-f)(x)|\leq\|g-f\|\cdot\|x\|\) 又说明范数拓扑不弱于弱-*\WeakStarJoin{}拓扑。这些事实都不要求 \(X\) 完备。

\ProseParagraphBreak

记号中保留 \(X\) 十分重要。这里不是任意给 \(X^*\) 降低一个拓扑，而是明确规定只用来自 \(X\) 的求值泛函作为测试。Fremlin 的《Measure Theory》\cite[Appendix to Volume 2, 2A5I]{Fremlin2016Measure}把弱拓扑与弱-*\WeakStarJoin{}拓扑放在同一框架下讨论，适合用来对照这两种测试族；下面则用序列空间把区别具体化。

\subsection{弱-*\WeakStarJoin{}收敛}
\label{b1:subsec:090202}

\begin{definition}\label{b1:def:weak-star-convergence}
设 \((f_n)\subset X^*\) 且 \(f\in X^*\)。若 \(f_n(x)\to f(x)\) 对每个 \(x\in X\) 都成立，则称它\term[ruo-xing-shoulian]{弱-*\WeakStarJoin{}收敛}[weak-* convergence]于 \(f\)，记为 \(f_n\overset{*}{\rightharpoonup}f\)。
\end{definition}
\symidx{\(f_n\overset{*}{\rightharpoonup}f\)：泛函序列的弱-*\WeakStarJoin{}收敛}

这就是对偶空间内部的逐点收敛。条件 \(f\in X^*\) 是定义的一部分，不能只从逐点极限公式中省略。事实上，在赋予一致范数的 \(c_{00}\) 上令
\[
 f_n(x)=\sum_{k=1}^n x_k.
\]
每个 \(f_n\) 都连续；对每个有限支撑的 \(x\)，它逐点趋于 \(\sum_{k\geq1}x_k\)，但这个极限泛函在一致范数下不连续，因为它在 \(x^{(N)}=\sum_{k=1}^Ne_k\) 上取值 \(N\)，而 \(\|x^{(N)}\|_\infty=1\)。所以当原空间不完备时，连续泛函的逐点极限未必仍是弱-*\WeakStarJoin{}拓扑中的点。

\ProseParagraphBreak

有限多个输入对应的起始指标取最大值，便得到序列在弱-*\WeakStarJoin{}拓扑中的收敛；网也有同样的逐点判别，只须把最大值换成有限个指标的共同上界。以下有界性命题则仍明确针对序列。

\begin{proposition}\label{b1:prop:weak-star-sequence-bounded}
若 \(X\) 是 Banach 空间，且 \(f_n\overset{*}{\rightharpoonup}f\) 于 \(X^*\)，则 \(\sup_n\|f_n\|<\infty\)。
\end{proposition}
\begin{proof}
对每个 \(x\in X\)，收敛的标量列 \(\bigl(f_n(x)\bigr)\) 有界。把 \(f_n\) 看成从 Banach 空间 \(X\) 到 \(\mathbb K\) 的有界线性算子，\cref{b1:thm:uniform-boundedness}直接给出结论。
\end{proof}

这里的一致有界原理作用在 \(X\) 上，而\cref{b1:prop:weak-sequence-bounded}中作用在自动完备的 \(X^*\) 上，故两条命题的假设不同。\subsecref{b1:subsec:080302}中的例子恰好说明完备性不能直接删去：在赋予一致范数的 \(c_{00}\) 上，令 \(f_n(x)=nx_n\)。每个固定输入只有有限个非零坐标，所以 \(f_n(x)\) 最终等于零，即 \(f_n\overset{*}{\rightharpoonup}0\)；但 \(\|f_n\|=n\)。对偶空间 \(c_{00}^*\) 本身完备，并不能补上定义域的这个缺口。

\ProseParagraphBreak

为了把弱-*\WeakStarJoin{}收敛写成常见的坐标条件，我们先确定一个具体对偶空间。

\begin{proposition}[\(\ell^1\) 的完备性与连续对偶]\label{b1:prop:l1-dual}
空间 \(\ell^1\) 关于范数 \(\|a\|_1=\sum_k|a_k|\) 完备。对 \(u\in\ell^\infty\)，公式
\[
 F_u(a)=\sum_{k=1}^{\infty}a_ku_k\quad(a\in\ell^1)
\]
定义 \((\ell^1)^*\) 中的元素。映射 \(u\mapsto F_u\) 是 \(\ell^\infty\) 到 \((\ell^1)^*\) 的线性等距满射。
\end{proposition}
\begin{proof}
先设 \((a^{(j)})\) 是 \(\ell^1\) 中的 Cauchy 列。每个坐标列都是 Cauchy 列，记 \(a_k=\lim_j a_k^{(j)}\)。给定 \(\varepsilon>0\)，取 \(J\) 使 \(j,m\geq J\) 时 \(\|a^{(j)}-a^{(m)}\|_1<\varepsilon\)。固定 \(j\geq J\)，对任意 \(N\) 令 \(m\to\infty\)，得到
\[
 \sum_{k=1}^N|a_k^{(j)}-a_k|\leq\varepsilon.
\]
令 \(N\to\infty\)，可知 \(a-a^{(j)}\in\ell^1\) 且 \(\|a-a^{(j)}\|_1\leq\varepsilon\)。因此 \(a\in\ell^1\) 且 \(a^{(j)}\to a\)，完备性得证。

绝对收敛性和估计 \(|F_u(a)|\leq\|u\|_\infty\cdot\|a\|_1\) 给出 \(\|F_u\|\leq\|u\|_\infty\)。逐个代入单位坐标向量 \(e_k\)，得到 \(\|F_u\|\geq|u_k|\)，取上确界即有范数相等。

反过来，给定 \(F\in(\ell^1)^*\)，令 \(u_k=F(e_k)\)，则 \(|u_k|\leq\|F\|\)，所以 \(u\in\ell^\infty\)。对 \(a\in\ell^1\)，其有限坐标截断 \(a^{(N)}\) 在 \(\ell^1\) 范数下趋于 \(a\)，故
\[
 F(a)=\lim_{N\to\infty}F(a^{(N)})
 =\lim_{N\to\infty}\sum_{k=1}^N a_ku_k=F_u(a).
\]
各坐标 \(u_k\) 已由 \(F(e_k)\) 唯一确定，因而表示唯一。
\end{proof}

上述配对在实、复情形下都不插入共轭；这样，\(a\mapsto F_u(a)\) 与 \(u\mapsto F_u\) 都是线性的。它不同于 Hilbert 空间的内积表示约定。

\begin{corollary}\label{b1:cor:linfty-weak-star-coordinate}
在识别 \(\ell^\infty=(\ell^1)^*\) 后，序列 \((u^{(n)})\) 弱-*\WeakStarJoin{}收敛于 \(u\)，当且仅当
\[
 \sup_n\|u^{(n)}\|_\infty<\infty,
 \quad u_k^{(n)}\longrightarrow u_k\quad\text{对每个 }k.
\]
\end{corollary}
\begin{proof}
弱-*\WeakStarJoin{}收敛时，用 \(e_k\in\ell^1\) 测试便得到坐标收敛。\Cref{b1:prop:l1-dual}已经直接证明 \(\ell^1\) 是 Banach 空间，故\cref{b1:prop:weak-star-sequence-bounded}给出统一界。

反过来，设上述统一界为 \(M\)。坐标极限满足 \(|u_k|\leq M\)，所以 \(u\in\ell^\infty\)。对任意 \(a\in\ell^1\)，分成前 \(N\) 项与尾部可得
\[
 \left|\sum_{k=1}^{\infty}a_k(u_k^{(n)}-u_k)\right|
 \leq\sum_{k=1}^N|a_k|\cdot|u_k^{(n)}-u_k|
      +2M\sum_{k>N}|a_k|.
\]
先选 \(N\) 使尾部任意小，再由有限个坐标的收敛使首项任意小，即得到每个 \(a\) 对应的求值收敛。
\end{proof}

有限个坐标测试本身看不到尾部；统一范数界与测试序列的绝对可和性一起，才把尾部误差控制住。这个先截断再取极限的步骤，也解释了为什么不能只凭坐标收敛就断言弱-*\WeakStarJoin{}收敛。

\subsection{弱拓扑与弱-*\WeakStarJoin{}拓扑的区别}
\label{b1:subsec:090203}

把 \(X^*\) 当作赋范空间，它自己的弱拓扑是 \(\sigma(X^*,X^{**})\)，允许使用所有 \(\Phi\in X^{**}\) 测试；弱-*\WeakStarJoin{}拓扑 \(\sigma(X^*,X)\) 只允许其中形如 \(J_Xx\) 的测试。因此在 \(X^*\) 上有
\[
 \text{范数拓扑}\ \supseteq\ \sigma(X^*,X^{**})
 \ \supseteq\ \sigma(X^*,X).
\]
这里的包含表示开集族的包含，相应地，范数收敛蕴含弱收敛，弱收敛蕴含弱-*\WeakStarJoin{}收敛。如果 \(J_X(X)=X^{**}\)，后两种拓扑便相同；一般情况下，缺少的测试确实可能改变收敛性。

\ProseParagraphBreak

自然嵌入本身还有一个拓扑性质，后面把原空间放进二次对偶时会反复使用。

\begin{proposition}\label{b1:prop:canonical-weak-embedding}
自然嵌入 \(J_X\) 把 \(\bigl(X,\sigma(X,X^*)\bigr)\) 同胚地映到 \(X^{**}\) 中的像 \(J_X(X)\)，其中像空间采用 \(X^{**}\) 上弱-*\WeakStarJoin{}拓扑 \(\sigma(X^{**},X^*)\) 的子空间拓扑。
\end{proposition}
\begin{proof}
这里要说明 \(J_X\) 是单射。而 \(X^{**}\) 中以 \(J_Xx\) 为中心、由有限个 \(f_1,\ldots,f_m\in X^*\) 给出的弱-*\WeakStarJoin{}基本邻域，在 \(J_X\) 下的反像由
\[
 |(J_Xy-J_Xx)(f_j)|=|f_j(y-x)|<\varepsilon
 \quad(1\leq j\leq m)
\]
刻画。这恰好是 \(X\) 的基本弱邻域，且所有基本弱邻域都如此得到。因此 \(J_X\) 及其在像上的逆映射都连续。
\end{proof}

来看一个弱-*\WeakStarJoin{}收敛却不弱收敛的序列。在实空间 \(\ell^\infty=(\ell^1)^*\) 中，令
\[
 u_k^{(n)}=\mathbf1_{\{k\geq n\}}.
\]
它的范数恒等于 \(1\)，每个固定坐标都最终变成零，所以由\cref{b1:cor:linfty-weak-star-coordinate}有 \(u^{(n)}\overset{*}{\rightharpoonup}0\)。另一方面，令 \(c\subset\ell^\infty\) 为实收敛序列组成的子空间。通常极限映射 \(\lambda:c\to\mathbb R\) 是范数为一的连续线性泛函，由\cref{b1:thm:hahn-banach}可将它保范延拓为某个 \(L\in(\ell^\infty)^*\)。对每个固定 \(n\)，序列 \((u_k^{(n)})_{k\geq1}\) 的通常极限为 \(1\)，因此 \(L(u^{(n)})=1\)。这个合法的弱测试不趋于零，故 \((u^{(n)})\) 不弱收敛于零。它也不可能弱收敛于别的向量，因为弱收敛蕴含弱-*\WeakStarJoin{}收敛，而弱-*\WeakStarJoin{}极限唯一。

\ProseParagraphBreak

这里使用尾部示性序列是有意的：单位坐标向量 \(e_n\) 在 \(\ell^\infty\) 中反而弱收敛于零。事实上，对任意 \(L\in(\ell^\infty)^*\)，在每个有限和中选择模为 \(1\) 的标量 \(\alpha_k\)，使 \(\alpha_k\cdot L(e_k)=|L(e_k)|\)，便有
\[
 \sum_{k=1}^N|L(e_k)|
 =L\left(\sum_{k=1}^N\alpha_ke_k\right)\leq\|L\|.
\]
故 \(L(e_n)\to0\)。相同的坐标外观并不足以判断弱收敛，必须看清所处的空间及它允许的全部泛函。

\ProseParagraphBreak

弱-*\WeakStarJoin{}拓扑保留逐点测试，却舍去了一部分弱测试。下一节将说明，这种舍弃换来一个重要的紧致性结论；使用那个结论时，既要记住原空间是谁，也不能把一般拓扑的紧致性预先当作子列收敛性。
